\section{Collision kernels and a velocity majorant}
\label{ker:section}\label{sec:kernel}

We establish gain bounds for narrow velocity distributions and for
the weight $H$. They control the finite evolutions and the global
continuation. The sheet and quadratic estimates are adapted from
\cite[Section~3]{OpenAIPeriodic}; their proofs are included below,
followed by the small-mass estimate needed for continuation.

Write $v=(\bar v,v_3)\in\R^2\times\R$ and set
\[
 W(v)=1+|\bar v|^{-q},\qquad
 H(v)=e^{-|v|^2}\langle v\rangle^{-m}W(v),
 \qquad 1<q<2,\quad m=100,
\]
as in \eqref{eq:weight}. Values on the velocity axis are irrelevant to
all the almost-everywhere estimates below. The function $H$ has every
polynomial moment and belongs to $L^p(\R^3)$ whenever $1\leq p<2/q$.
Indeed, its only local singularity is $|\bar v|^{-q}$ in two dimensions,
and its decay at infinity is Gaussian. Every estimate below is proved
for an arbitrary fixed $q\in(1,2)$; it may therefore also be applied
with a different fixed exponent in that interval.

A model for the cold velocities, in parameters
$w=(\bar w,w_3)\in\R^3$, is
\[
 \bar p=d\bar w,\qquad p_3=\upsilon(z)+dw_3,
\]
where $d>0$ is small and $\upsilon$ is bounded, nondecreasing, and
$C^1$. At a fixed position after free transport, the axial parameter
relation and its Jacobian become
\[
 p_3=\upsilon(z-tp_3)+dw_3,\qquad
 \frac{\partial p_3}{\partial w_3}
   =\frac{d}{1+t\upsilon'(z-tp_3)}.
\]
A fixed parameter law can therefore produce strong axial
concentration. We will estimate gains using transverse Gaussian
majorants while imposing no density bound on the axial law. For any
nonnegative velocity density $c$ and $Z\geq0$, positivity gives
\[
 0\leq h\leq ZH\quad\Longrightarrow\quad
 \frac{P(c,h)}{H}\leq Z\frac{P(c,H)}{H}.
\]
Thus $P(c,H)/H$ controls the gain in the weighted supremum. The
collision formula below realizes this coefficient as the mass of a
normalized measure on the two input velocities.

\subsection{The exact Carleman formula}

The hyperplane representation used here is a Carleman formula;
see also \cite[Appendix~C]{GambaPanferovVillani2009}.
We derive its coefficient for the normalization of this paper.

\begin{lemma}[Collision geometry]\label{lem:carleman}
For the normalization in \eqref{eq:collision}, the symmetric bilinear
gain is
\begin{equation}\label{ker:carleman}
 P(c,h)(v)
 =2\int_{\R^3}c(p)\,\dd p
       \int_{\R^3}
       \delta\bigl((v-p)\cdot(\xi-v)\bigr)h(\xi)\,\dd\xi.
\end{equation}
For fixed $p\ne v$, the inner delta denotes surface measure on the
plane $(v-p)\cdot(\xi-v)=0$, divided by $|v-p|$.
For incoming velocities $p,\xi$, the outgoing measure has total mass
$2\pi|p-\xi|$ and is uniform on the sphere with diameter $[p,\xi]$.
Formula~\eqref{ker:carleman} also defines the gain when $c$ is a
nonnegative velocity measure and $h$ is a nonnegative density, whenever
the displayed integral is finite almost everywhere.
\end{lemma}

\begin{proof}
Put $g=p-\xi$. Reflection sends $g$ to
$g-2(g\cdot\omega)\omega$. On either hemisphere of the normal sphere,
the change from the polar angle of $\omega$ to the reflected direction
shows that $|g\cdot\omega|\,\dd\omega$ pushes forward to
$(|g|/4)\,\dd\sigma$. The two hemispheres therefore give
$(|g|/2)\,\dd\sigma$, where $\sigma\in S^2$ is the outgoing relative
direction. Equivalently, for every nonnegative test function $\psi$,
the one-output measure is
\begin{equation}\label{ker:sphere-law}
 \frac{|p-\xi|}{2}\int_{S^2}
 \psi\left(\frac{p+\xi}{2}
                 +\frac{|p-\xi|}{2}\sigma\right)\dd\sigma.
\end{equation}
Its mass is $2\pi|p-\xi|$. The affine map in this integral multiplies
surface area by $|p-\xi|^2/4$, so the density on its image sphere is
$2/|p-\xi|$.

The same sphere is the zero set of
$f(v)=(v-p)\cdot(\xi-v)$. On that sphere,
$|\nabla_v f(v)|=|p-\xi|$. The coarea formula consequently identifies
the measure in \eqref{ker:sphere-law} with $2\delta(f(v))\,\dd v$.
Integration against the two incoming densities and Tonelli's theorem
give \eqref{ker:carleman}. The usual pre/post-collisional change of
variables identifies this incoming-pair expression with the gain in
\eqref{eq:collision}. Coincident incoming velocities have zero
collision rate. The same identities can first be integrated against a
velocity measure $c(\dd p)$, proving the asserted extension.
\end{proof}

In the cutoff gain, the integrand in \eqref{ker:carleman} is multiplied
by $\min(1,b/|p-\xi|)$. The relative input speed equals the relative
output speed, so this factor preserves the collision symmetries.
It is at most one. All upper estimates below consequently hold for
$P^b$ as well, uniformly in $b\geq1$.

\subsection{A geometric slice estimate}

For $d>0$, let $\mathcal G_d$ be the centered transverse Gaussian
probability measure with density
\[
 \frac{1}{\pi d^2}e^{-|\bar p|^2/d^2},
\]
and let $\mathcal G_0=\delta_0$. We shall also use the uniform
probability measure $\mathcal U_d$ on $\{\bar p:|\bar p|<d\}$.

\begin{lemma}[Transverse averaging and axial slices]
\label{ker:geometric-slice}
Let $\lambda$ be a probability mixture of the measures
$\mathcal U_d$, $d>0$, and $\delta_0$, and let $\rho$ be any axial
probability measure. No bound on its density or its support is
required. Put $c=\lambda\otimes\rho$. For $R\geq1$, $\bar v\ne0$,
and an interval $I\subset\R$ of length $0<\ell\leq2R$, one has
\begin{align}
 &\frac{1}{W(v)}\int c(\dd p)
   \int_{|\xi-v|\leq R}
    \delta\bigl((v-p)\cdot(\xi-v)\bigr)W(\xi)\,\dd\xi
     \leq C_q R^2,\label{ker:geometric-total}\\
 &\frac{1}{W(v)}\int c(\dd p)
   \int_{\substack{|\xi-v|\leq R\\ \xi_3\in I}}
    \delta\bigl((v-p)\cdot(\xi-v)\bigr)W(\xi)\,\dd\xi
     \leq C_q\bigl(R\ell+R^{q-1}\ell^{2-q}\bigr).
     \label{ker:geometric-interval}
\end{align}
The constants are independent of all transverse widths and of $\rho$.
In particular, these statements hold with $\lambda=\mathcal G_d$ for
every $d\geq0$.
\end{lemma}

\begin{proof}
Set $r=|\bar v|$, $n=\bar v-\bar p$, and $e_n=n/|n|$. Except on a
set of transverse measure zero, $n\ne0$. For fixed $p_3,\xi_3$, the
delta in the transverse variables is line measure divided by $|n|$,
on a line whose signed distance from zero is
\[
 A-B(\xi_3-v_3),\qquad
 A=e_n\cdot\bar v,\quad B=\frac{v_3-p_3}{|n|}.
\]
The regular part of its integral over $|\xi-v|\leq R$ is at most
$2R/|n|$. The singular part is bounded by its integral over the full
line, namely
\begin{equation}\label{ker:line-integral}
 \frac{1}{|n|}\int_{\R}
      \bigl((A-B(\xi_3-v_3))^2+s^2\bigr)^{-q/2}\,\dd s
 =\frac{a_q}{|n|}|A-B(\xi_3-v_3)|^{1-q},
\end{equation}
where $a_q=\int_{\R}(1+s^2)^{-q/2}\,\dd s<\infty$ because $q>1$.
Infinite values on a zero-distance slice are harmless in the
nonnegative iterated integral.

For every interval $J\subset[-R,R]$ of length at most $\ell$,
\begin{equation}\label{ker:scalar-interval}
 \int_J |A-Bx|^{1-q}\,\dd x
 \leq C_q R^{q-1}\ell^{2-q}|A|^{1-q}.
\end{equation}
If $2R|B|\leq|A|$, the integrand is at most
$C_q|A|^{1-q}$, and the assertion follows from $\ell\leq2R$.
Otherwise, integration about the possible root gives an upper bound
$C_q|B|^{1-q}\ell^{2-q}$; since $1-q<0$ and
$|B|>|A|/(2R)$, this proves \eqref{ker:scalar-interval}. When $A=0$
the displayed upper bound is understood in the extended sense.

We next prove
\begin{equation}\label{ker:transverse-expectation}
 \int\frac{1+|e_n\cdot e_{\bar v}|^{1-q}}{|n|}\,
          \lambda(\dd\bar p)\leq\frac{C_q}{r}.
\end{equation}
It suffices to consider a disk of radius $d$. If $d\leq r/2$, then
$|n|\geq r/2$ and $e_n\cdot e_{\bar v}\geq1/3$, which proves the
bound. If $d>r/2$, use polar coordinates for $n$. Its range lies in
$|n|<r+d<3d$. The factor $|n|^{-1}$ cancels the radial Jacobian, while
\[
 \int_0^{2\pi}\bigl(1+|\cos\theta|^{1-q}\bigr)\,\dd\theta<\infty
\]
because $q<2$. Division by the disk area gives $C_q/d\leq2C_q/r$.
The point mass at zero satisfies the same estimate directly.
Finally, the Gaussian is a probability mixture of these disks: for
$d>0$ its mixture density in the disk radius $a$ is
\[
 \frac{2a^3}{d^4}e^{-a^2/d^2}\,\dd a,
\]
whose integral is one. This proves the Gaussian assertion as well.

Integrating \eqref{ker:line-integral} in $\xi_3$, applying
\eqref{ker:scalar-interval}, and then
\eqref{ker:transverse-expectation}, bounds the expression in
\eqref{ker:geometric-interval} by
\[
 C_q R\ell\frac{r^{-1}}{1+r^{-q}}
 +C_qR^{q-1}\ell^{2-q}\frac{r^{-q}}{1+r^{-q}}.
\]
Both fractions are bounded for $q>1$. These bounds are independent of
$p_3$, so integration against the arbitrary probability $\rho$
preserves them. Taking the entire axial range of length $2R$ also
gives \eqref{ker:geometric-total}.
\end{proof}

\subsection{Sheets and normalized kernel measures}

\begin{definition}[Uniform sheet class]\label{ker:sheet-class}
Fix $L\geq1$. An elementary sheet in $\mathcal S_L$ is a product
probability measure
\begin{equation}\label{ker:sheet-definition}
 c(\dd p)=\mathcal G_d(\dd\bar p)\rho(\dd p_3),\qquad
 0\leq d\leq\tfrac18,\qquad
 \int_{\R}e^{8a^2}\rho(\dd a)\leq L.
\end{equation}
We also allow probability mixtures of elementary sheets. An upper
bound by a sheet with amplitude $A_0$ means domination as measures by
$A_0$ times such a mixture.
\end{definition}

Every axial probability supported in $[-A,A]$ satisfies this
definition with $L=e^{8A^2}$. In particular, axial atoms and limits of
arbitrarily narrow axial densities are permitted. The exponential
moment is used only to remove large axial velocities, and imposes no
axial regularity.

For such a sheet and $\bar v\ne0$, define its normalized kernel by
\begin{equation}\label{ker:normalized-kernel}
 \mathcal K_c(v;\dd p\,\dd\xi)
 =2\,c(\dd p)\,
       \delta\bigl((v-p)\cdot(\xi-v)\bigr)
       \frac{H(\xi)}{H(v)}\,\dd\xi.
\end{equation}
The total mass is $P(c,H)(v)/H(v)$. Kernel restrictions always refer
to this complete measure, including the ratio $H(\xi)/H(v)$.

\begin{lemma}[Uniform sheet gain bound]\label{lem:sheet-bound}
For every $L<\infty$ there is $C_{q,m,L}<\infty$ such that
\begin{equation}\label{ker:sheet-gain}
 P(c,H)(v)\leq C_{q,m,L}H(v)
 \quad\hbox{for almost every }v,
 \qquad c\in\mathcal S_L.
\end{equation}
More quantitatively, for $R\geq1$,
\begin{equation}\label{ker:shell-tail}
 \mathcal K_c\bigl(v;
      R<|p|+|\xi-v|\leq2R\bigr)
 \leq C_{q,m,L}R^2e^{-cR^2},
\end{equation}
where $c>0$ is an absolute constant. The same conclusions, multiplied
by $A_0$, hold under sheet domination with amplitude $A_0$.
\end{lemma}

\begin{proof}
By positivity it suffices to treat an elementary sheet. Put
$k=\xi-v$ and $w=p+k$. On the collision plane,
\begin{equation}\label{ker:energy-factor}
 |p|^2+|\xi|^2=|v|^2+|w|^2,
 \qquad
 \frac{H(\xi)}{H(v)}
 =e^{|p|^2-|w|^2}
    \frac{\langle v\rangle^m}{\langle\xi\rangle^m}
    \frac{W(\xi)}{W(v)}.
\end{equation}
The same identity implies
$\langle v\rangle\leq\langle p\rangle\langle\xi\rangle$.
Since $\langle p\rangle^m\leq C_m e^{|p|^2}$, the kernel is bounded
by
\begin{equation}\label{ker:tilted-kernel}
 C_m e^{2|p|^2}c(\dd p)\,
 e^{-|w|^2}\delta\bigl((v-p)\cdot k\bigr)
       \frac{W(v+k)}{W(v)}\,\dd k.
\end{equation}

The tilted cold measure retains a product majorant. For $0<d\leq1/8$,
\begin{equation}\label{ker:transverse-tilt}
 e^{2|\bar p|^2}\mathcal G_d(\dd\bar p)
 \leq4e^{-46|\bar p|^2}\mathcal G_{2d}(\dd\bar p).
\end{equation}
This follows by dividing the two Gaussian densities: their ratio is
$4\exp[-(3/(4d^2)-2)|\bar p|^2]$, and
$3/(4d^2)-2\geq46$. The assertion at $d=0$ has the evident
interpretation. The axial tilted measure
$\widetilde\rho(\dd a)=e^{2a^2}\rho(\dd a)$ has mass at most $L$,
and satisfies
\begin{equation}\label{ker:axial-tilt-tail}
 \widetilde\rho\{|a|>u\}\leq Le^{-6u^2},\qquad u\geq0.
\end{equation}
Consequently $e^{2|p|^2}c(\dd p)$ is bounded by $4L$ times a product
of a centered transverse Gaussian and an axial probability. Moreover,
its restriction to $|p|>u$ is bounded by
$CLe^{-cu^2}$ times a sum of at most two such product probabilities.
To see the last assertion, split into $|\bar p|>u/\sqrt2$ and
$|p_3|>u/\sqrt2$, and apply
\eqref{ker:transverse-tilt} and \eqref{ker:axial-tilt-tail},
respectively. The axial probabilities obtained by normalization may
have large support; Lemma~\ref{ker:geometric-slice} permits this.

On the shell in \eqref{ker:shell-tail}, either $|p|>R/4$ or
$|w|>R/2$, because $|p|+|k|\leq2|p|+|w|$. In the first case the
preceding tilted-measure tail contributes $CLe^{-cR^2}$. In the
second case $e^{-|w|^2}\leq e^{-R^2/4}$. In either case the remaining
integral has $|k|\leq2R$ and is bounded by $C_qR^2$ using
Lemma~\ref{ker:geometric-slice}. This proves
\eqref{ker:shell-tail}. The base region $|p|+|k|\leq2$ is bounded
by the same geometric lemma and the total tilted mass. Summing the
dyadic shells proves \eqref{ker:sheet-gain} and also the finiteness
of \eqref{ker:normalized-kernel}.
\end{proof}

The sheet gain bound controls the total normalized mass. We also
need uniform control of the input tails and decay for large output
speed, so that later path estimates can be restricted to bounded
velocity regions.

\begin{lemma}[Uniform tightness and large-output decay]
\label{lem:kernel-tightness}
Fix $1\leq L<\infty$. Uniformly for $c\in\mathcal S_L$ and outgoing
velocities off the axis, the tails in $|p|+|\xi-v|$ are negligible:
\begin{equation}\label{ker:uniform-tail}
 \lim_{R\to\infty}\sup_{c\in\mathcal S_L}\esssup_v
 \mathcal K_c(v;|p|+|\xi-v|>R)=0.
\end{equation}
Moreover,
\begin{equation}\label{ker:large-output}
 \lim_{M\to\infty}\sup_{c\in\mathcal S_L}
 \esssup_{|v|>M}\mathcal K_c(v;\R^3\times\R^3)=0.
\end{equation}
Both conclusions hold under uniformly bounded sheet amplitudes.
\end{lemma}

\begin{proof}
The dyadic estimate \eqref{ker:shell-tail} proves
\eqref{ker:uniform-tail}. Write $k=\xi-v$. We may consequently work on
$|p|+|k|\leq R$ and remove that restriction at the end of each
argument. On this region the exponential and polynomial factors in
\eqref{ker:energy-factor} are bounded by a constant depending on $R$.

For \eqref{ker:large-output}, if $|\bar v|>2R+1$, then
$|\bar v-\bar p|\geq|\bar v|/2$ and
$|\bar\xi|\geq|\bar v|-R$. The transverse line estimate, followed
by the axial integration over a length at most $2R$, gives
$C_R/|\bar v|$. If $|v_3|>2R+1$, integrate the delta in $\xi_3$
instead. Its divisor is $|v_3-p_3|\geq|v_3|/2$, and
\begin{equation}\label{ker:disk-singularity}
 \sup_{a\in\R^2}\int_{|\zeta-a|\leq R}
       (1+|\zeta|^{-q})\,\dd\zeta
 \leq C_q(R^2+R^{2-q}).
\end{equation}
For completeness, if $|a|\leq2R$ the disk is contained in the disk
of radius $3R$ about zero; if $|a|>2R$, then $|\zeta|>|a|/2>R$ on
the disk, giving the same bound. We obtain $C_R/|v_3|$. At least one
of the two outgoing components becomes large with $|v|$, proving
\eqref{ker:large-output}.
\end{proof}

\subsection{Gaussian parameter laws and compressed axial maps}

The next formulation records the uniformity needed when a cold
profile is parameterized before its axial change of variables.

\begin{lemma}[Gaussian parameter tails]
\label{ker:parameter-sheets}
Let a measure in parameters $w=(\bar w,w_3)\in\R^2\times\R$ have
density $j$ satisfying
\[
 0\leq j(w)\leq C_0\pi^{-3/2}e^{-|w|^2}.
\]
Let its velocity image be given by
\begin{equation}\label{ker:parameter-map}
 \bar p=d\bar w,\qquad p_3=\Phi(w_3),\qquad
 0\leq d\leq\tfrac1{16},\qquad
 |\Phi(a)|\leq A+\delta|a|,\quad0\leq\delta\leq\tfrac1{16},
\end{equation}
where $\Phi$ is merely measurable. The image measure is dominated
by $C_0$ times a sheet in $\mathcal S_{L_A}$, with $L_A$ depending
only on $A$. Its gain satisfies Lemma~\ref{lem:sheet-bound} with a
constant depending only on $q,m,A,C_0$.

If $c_{>B}$ is the image of the restricted parameter measure
$\ind_{|w|>B}j(w)\,\dd w$, then
\begin{equation}\label{ker:parameter-tail}
 \esssup_v\frac{P(c_{>B},H)(v)}{H(v)}
 \leq C_{q,m,A,C_0}e^{-B^2/2},\qquad B\geq0.
\end{equation}
Thus parameter truncations are valid in the normalized collision
measure, even when $\Phi$ compresses intervals or has atoms in its
image law. Probability mixtures of measurable families with common
$q,m,A,C_0$ and the width bounds in \eqref{ker:parameter-map} preserve
these assertions. More generally, a positive mixing measure of total
mass $B_0<\infty$ multiplies the domination amplitude and quantitative
upper bounds by $B_0$; the sheet parameter $L_A$ is unchanged.
For \eqref{ker:parameter-tail}, impose the restriction $|w|>B$ in
each component before mixing. Additional measurable masks taking
values in $[0,1]$ preserve the assertions.
\end{lemma}

\begin{proof}
The Gaussian majorant is a product measure. Its transverse image is
$\mathcal G_d$, and its axial image is the probability
$\rho=\Phi_*(\pi^{-1/2}e^{-a^2}\,\dd a)$. There is no need to
differentiate or invert $\Phi$. From \eqref{ker:parameter-map},
\[
 e^{8|\Phi(a)|^2}\leq e^{16A^2}e^{16\delta^2a^2}.
\]
Even under the wider axial probability
$(2\pi)^{-1/2}e^{-a^2/2}\,\dd a$, the integral of the right side
is
\[
 e^{16A^2}(1-32\delta^2)^{-1/2}\leq2e^{16A^2}.
\]
We may therefore take $L_A=2e^{16A^2}$ for both the original and
the wider parameter laws. The original image has the required
sheet domination.

On the parameter tail,
\[
 \ind_{|w|>B}e^{-|w|^2}
 \leq e^{-B^2/2}e^{-|w|^2/2}.
\]
After normalization the right side is $2^{3/2}e^{-B^2/2}$ times a
Gaussian probability with parameter widths enlarged by $\sqrt2$.
Its transverse velocity width is at most $\sqrt2/16<1/8$, and
its axial image has the exponential moment just proved.
Lemma~\ref{lem:sheet-bound} applied to this dominating image gives
\eqref{ker:parameter-tail}.

For a probability mixture, mix the dominating sheets just constructed;
this is again a member of $\mathcal S_{L_A}$. If the mixing measure
has mass $B_0>0$, normalize it by $B_0$ and retain that factor in
the domination amplitude. If $B_0=0$, the mixed image is zero.
The same argument applied to the wider Gaussian majorants of the
component tails proves the asserted factor in
\eqref{ker:parameter-tail}. Thus the estimates are uniform over
families of mixtures whose masses $B_0$ have a common finite bound
independent of the tail cutoff. Masks taking values in $[0,1]$ only
decrease the measures, whether imposed before or after the pushforward.
\end{proof}

Return to the profile at the start of the section:
$p_3=\upsilon(z-tp_3)+d w_3$, with $\upsilon$ nondecreasing, and
let $u$ solve $u=\upsilon(z-tu)$. The function
$p\mapsto p-\upsilon(z-tp)$ is increasing with inverse Lipschitz
constant at most one. Hence
\begin{equation}\label{ker:compressed-displacement}
 |p_3-u|\leq d|w_3|.
\end{equation}
If $|u|\leq A$, Lemma~\ref{ker:parameter-sheets} applies. A large
Jacobian compression of the axial law does not affect its constants.
Gaussian majorants with any fixed positive exponent in place of one
are covered by a fixed rescaling of the parameters and a corresponding
decrease of the upper bounds on the physical widths.

\subsection{The hot--hot gain}

\begin{lemma}[Quadratic gain bound]\label{lem:hot-bound}
There is a constant $C_{q,m}$ such that
\begin{equation}\label{ker:hot-gain}
 P(H,H)\leq C_{q,m}H
 \quad\hbox{almost everywhere}.
\end{equation}
In fact the proof below works for every $m>5$.
\end{lemma}

\begin{proof}
The integrand is symmetric in $p,\xi$. On the collision plane,
$|v|^2\leq|p|^2+|\xi|^2$, so at least one incoming speed is at
least $|v|/2$. It suffices, at a factor of two, to retain
$|\xi|\geq|v|/2$. Formula~\eqref{ker:energy-factor} now gives
\begin{align}
 \frac{P(H,H)(v)}{H(v)}
 &\leq C_m\int_{|\xi|\geq|v|/2}
 \delta\bigl((v-p)\cdot(\xi-v)\bigr)
 e^{-|p+\xi-v|^2}\langle p\rangle^{-m}
 \frac{W(p)W(\xi)}{W(v)}\,\dd p\,\dd\xi,
 \label{ker:hot-reduced}
\end{align}
because $\langle v\rangle^m/\langle\xi\rangle^m\leq2^m$.

Again write $k=\xi-v$ and $w=p+k$. On
$R<|p|+|k|\leq2R$, $R\geq1$, either $|p|>R/4$ or $|w|>R/2$.
Consequently
\begin{equation}\label{ker:hot-shell-weight}
 \langle p\rangle^{-m}e^{-|w|^2}
 \leq C_mR^{-m}e^{-|w|^2/2}.
\end{equation}
In the second alternative this uses the finite supremum of
$s^m e^{-s^2/2}$; in the first it follows directly from the
polynomial factor.

The measure $W(p)\ind_{|p|\leq2R}\,\dd p$ is dominated by
\[
 \bigl(1+|\bar p|^{-q}\bigr)
 \ind_{|\bar p|\leq2R}\,\dd\bar p
 \ \otimes\ \ind_{|p_3|\leq2R}\,\dd p_3.
\]
Its transverse factor is a positive mixture of centered disks.
Indeed, for $D>0$,
\[
 |\bar p|^{-q}\ind_{|\bar p|\leq D}
 =D^{-q}\ind_{|\bar p|\leq D}
   +q\int_0^D a^{-q-1}\ind_{|\bar p|\leq a}\,\dd a
\]
almost everywhere. Upon normalizing the disks, the mixture has
total mass
$\pi D^2+2\pi D^{2-q}/(2-q)$. The axial factor has mass $4R$.
For $D=2R$ the resulting product mixture therefore has total mass
at most $C_qR^3$.

Discard $e^{-|w|^2/2}$ and the incoming-speed restriction in
\eqref{ker:hot-reduced}. Apply
Lemma~\ref{ker:geometric-slice} to each normalized product in this
mixture, with $|k|\leq2R$. In conjunction with
\eqref{ker:hot-shell-weight}, this yields the shell estimate
\begin{equation}\label{ker:hot-shell}
 \frac{P(H,H)(v)}{H(v)}\bigg|_{R<|p|+|\xi-v|\leq2R}
 \leq C_{q,m}R^{5-m}.
\end{equation}
The notation on the left denotes the restricted nonnegative
integral from \eqref{ker:hot-reduced}, including its fixed symmetry
factor. The base region is bounded by the same disk-mixture
argument with $R=1$. Summation over dyadic $R$ proves the result
for $m>5$.
\end{proof}

\subsection{A small-mass gain estimate}

The quadratic estimate controls a gain by a weighted supremum. For
global continuation we need a stronger conclusion when one input has
small mass while staying below a fixed multiple of $H$. The relevant
compactness is in the velocity integrals themselves: large speeds,
small transverse speeds, and the remaining collision diagonal can
all be removed with uniformly small normalized cost.

\begin{lemma}[Gain estimate from small mass]
\label{lem:gain-small-mass}
For each $0\leq B<\infty$ and $\varepsilon>0$ there is a finite
constant $C_{B,\varepsilon}$ such that every measurable density $j$
with $0\leq j\leq BH$ satisfies
\begin{equation}\label{ker:gain-small-mass}
 P(H,j)(v)\leq H(v)
 \left(\varepsilon+C_{B,\varepsilon}
                         \int_{\R^3}j(p)\,\dd p\right)
 \qquad\hbox{for almost every }v.
\end{equation}
The same constant works with $P^b$ in place of $P$.
\end{lemma}

\begin{proof}
The case $B=0$ is immediate. Assume $B>0$. We first establish
\begin{equation}\label{ker:hot-large-output}
 \lim_{M\to\infty}\esssup_{|v|>M}
                  \frac{P(H,H)(v)}{H(v)}=0.
\end{equation}
As in the proof of Lemma~\ref{lem:hot-bound}, symmetry reduces the
integral to $|\xi|\geq|v|/2$, up to a factor of two.
The shell estimate \eqref{ker:hot-shell} then removes
$|p|+|\xi-v|>R$ with an error tending to zero uniformly in $v$ as
$R\to\infty$. On the retained region, the right side of
\eqref{ker:hot-reduced} is bounded by a constant times
\[
 \int_{|p|\leq R}W(p)
 \left[\int_{|\xi-v|\leq R}
   \delta\bigl((v-p)\cdot(\xi-v)\bigr)
               \frac{W(\xi)}{W(v)}\,\dd\xi\right]\dd p.
\]
If $|\bar v|>2R+1$, the transverse divisor is at least
$|\bar v|/2$, and $W(\xi)$ is bounded on the retained region. The
transverse line integral and the axial interval of length $2R$ give
an upper bound $C_R/|\bar v|$ after integration in $p$. If
$|v_3|>2R+1$, integrating the delta in $\xi_3$ instead and using
\eqref{ker:disk-singularity} gives $C_R/|v_3|$. Here
$\int_{|p|\leq R}W(p)\,\dd p<\infty$, since $q<2$.
At least one output component tends to infinity with $|v|$.
First letting $|v|\to\infty$ and then $R\to\infty$ proves
\eqref{ker:hot-large-output}.

Choose $M$ so large that $BP(H,H)\leq\varepsilon H$ on $|v|>M$.
It remains to prove the desired estimate for $|v|\leq M$. On this
output ball, the part of the normalized quadratic gain with
$|p|+|\xi|>L$ tends uniformly to zero as $L\to\infty$. Indeed, this
deletion is symmetric in the two inputs. After the same reduction to
$|\xi|\geq|v|/2$, it implies
\[
 |p|+|\xi-v|>L-M,
\]
so \eqref{ker:hot-shell} bounds its cost by a summable tail.
The condition $j\leq BH$ transfers this estimate to $P(H,j)$.

Next retain both inputs in $|p|,|\xi|\leq L$. The Gaussian and
polynomial factors in the normalized quadratic integrand are then
bounded by a constant depending on $M,L$, leaving
\[
 \delta\bigl((v-p)\cdot(\xi-v)\bigr)
                       \frac{W(p)W(\xi)}{W(v)}.
\]
For $0<\delta<1$, the restriction $|\bar p|<\delta$ is dominated in
its $p$ variable by the product measure
\[
 (1+|\bar p|^{-q})\ind_{|\bar p|<\delta}\,\dd\bar p
       \otimes\ind_{|p_3|\leq L}\,\dd p_3.
\]
Its transverse factor is a positive mixture of centered disks, by
the layer-cake formula used in Lemma~\ref{lem:hot-bound}, and its
total mass is
\[
 2L\left(\pi\delta^2+
                  \frac{2\pi}{2-q}\delta^{2-q}\right).
\]
After normalization, \eqref{ker:geometric-total}, with radius at
least $L+M$, therefore bounds the cost of this restriction by
$C_{M,L}(\delta^2+\delta^{2-q})$. Symmetry gives the same bound
for $|\bar\xi|<\delta$. These estimates are uniform even when the
output velocity approaches its axis. Again $j\leq BH$ transfers
them to the mixed gain.

We have proved that, for a suitable compact set
\[
 K=\{p:|p|\leq L,\ |\bar p|\geq\delta\},
\]
the truncations $H_K=H\ind_K$ and $j_K=j\ind_K$ satisfy
\[
 0\leq P(H,j)-P(H_K,j_K)\leq\tfrac12\varepsilon H
       \qquad (|v|\leq M),
\]
uniformly over $0\leq j\leq BH$.
The function $H_K$ is bounded and supported in the ball of radius
$L$. By symmetry of $P$, put $j_K$ outside the plane integral in
\eqref{ker:carleman}. The area of a plane section of that ball is
at most $\pi L^2$, so
\[
 P(H_K,j_K)(v)\leq C_{L,\delta}
              \int_K\frac{j_K(p)}{|v-p|}\,\dd p.
\]
For $0<\rho<1$, the part with $|v-p|<\rho$ is at most
$C_{B,L,\delta}\rho^2$, because $j_K$ is bounded. The complement
is at most
$C_{L,\delta}\rho^{-1}\int j$.
Finally, $H$ has a positive lower bound on $|v|\leq M$ off its
null axis. Choose $\rho$ small enough that the near part is at most
$\varepsilon H/2$ there. Absorbing the other fixed factors into
$C_{B,\varepsilon}$ proves \eqref{ker:gain-small-mass}.
The cutoff assertion follows from $P^b\leq P$ on nonnegative inputs.
\end{proof}
